\chapter{Corrected transport, information and quantum contracts for tropical quivers}
\label{ch:toric-transport-audit}
\reviewstatus{This chapter adds explicit corrections to selected claims in the pinned Tropical Quivers and ToricGT research manuscripts. Seventeen finite declarations in TransportContracts.lean are kernel checked; the information-theoretic and stochastic-limit discussions are mathematical arguments with stated hypotheses, not a formalization of all source papers. None of the statements below extends the external forced Navier--Stokes theorem to the unforced problem.}

\section{Source scope and a quantitative contract}

The source revisions are:
\begin{itemize}
\item Tropical\_Quivers\_of\_Archs \cite{ToricAuditQuivers2026}, commit\\
\texttt{2b6ab79a685e2c6233fb8892e59589951feeddad};
\item ToricGT-Program, commit\\
\texttt{f91ba1811b5a4f70d53c7120acf8ca2a0049f85b}.
\end{itemize}
The pinned primary sources are available online.\footnote{Primary sources:
\url{https://github.com/amelie-iska/Tropical_Quivers_of_Archs/tree/2b6ab79a685e2c6233fb8892e59589951feeddad}
and \url{https://github.com/amelie-iska/ToricGT-Program/tree/f91ba1811b5a4f70d53c7120acf8ca2a0049f85b}.}
The companion review gives
the exact files, inspected sections, and unresolved coverage. The first
repository contains a manuscript and its PDF. The second contains five
research-program papers, earlier versions, reference PDFs, and a ToricGT
gitlink. The gitlink points at a separate code revision. A proposed
architecture, a numerical experiment and a proved theorem remain different
kinds of evidence.

The useful common interface is a finite typed representation with a
specified metric, update map, certificate and error budget. A surface
cellulation provides incidence maps; a gentle algebra additionally imposes
path relations; a physical fluid requires constitutive laws and analytic
estimates; a quantum channel requires complete positivity and trace
preservation. Each additional structure has its own contract. The following
corrections make those contracts compositional without identifying them.

\section{Suffix transport is the correct cyclic coboundary}

Let $V_0,\ldots,V_k$ be vector spaces, let
$A_j:V_{j-1}\to V_j$ be linear maps, and let $x_j\in V_j$. Define the edge
residual $r_j=x_j-A_jx_{j-1}$ and the suffix transport
\[
 T_j=A_kA_{k-1}\cdots A_{j+1}:V_j\longrightarrow V_k,
 \qquad T_k=I_{V_k}.
\]
Composition acts right to left. In particular $T_j$ transports a residual
from the target of edge $j$ to the endpoint of the path. No commutativity,
inverse, equal dimension or flatness assumption is required.

\begin{theorem}[Typed transport telescope]
\label{thm:typed-transport-telescope}
With the preceding notation,
\[
 \sum_{j=1}^{k}T_jr_j=x_k-A_k\cdots A_1x_0.
\]
For a closed path, $V_k=V_0$ and $x_k=x_0$, this equals $(I-H)x_0$, where
$H=A_k\cdots A_1$ is the based holonomy. It vanishes for every initial
vector if and only if $H=I$.
\end{theorem}
\begin{proof}
Expand the left side as
\[
 \sum_{j=1}^{k}A_k\cdots A_{j+1}x_j
 -\sum_{j=1}^{k}A_k\cdots A_jx_{j-1}.
\]
The term in the first sum with index $j<k$ equals the term in the second
sum with index $j+1$. These cancel in their common target space. The two
remaining terms are $x_k$ and $-A_k\cdots A_1x_0$. On a closed path, the
residual is zero on every vector precisely when the linear map $I-H$ is
zero. For a particular vector it is enough that $Hx_0=x_0$.
\end{proof}

The equal-space specialization is formalized by induction in
\path{TransportContracts.suffix_transport_telescope}. Its recursive
form is
\[
 s_0=0,\qquad s_{j+1}=A_{j+1}s_j+r_{j+1}.
\]
This is also an efficient implementation: the accumulated discrepancy is
transported once per edge, rather than storing all suffix matrices.

\begin{proposition}[A flat-cycle counterexample to prefix transport]
\label{prop:transport-prefix-counterexample}
Take scalar arrows $A_1=2$, $A_2=1/2$ and states
$x_0=x_2=0$, $x_1=1$. Then $A_2A_1=1$. Nevertheless the alternative
prefix-weighted expression
\[
 A_1r_1+(A_1A_2)r_2=\frac32
\]
is nonzero. The suffix expression $A_2r_1+r_2$ is zero.
\end{proposition}
\begin{proof}
Here $r_1=1$ and $r_2=-1/2$. Substitution gives
$2-1/2=3/2$ for prefixes and $1/2-1/2=0$ for suffixes.
\end{proof}

The pinned Tropical Quivers section on elliptic complexes used forward
prefix products as transport back to a cycle basepoint. The counterexample
shows that its asserted flatness implication does not follow from that
formula. Replacing prefixes by the suffixes above repairs the direction
and multiplication order simultaneously. The source is retained unchanged;
this chapter is an explicit amendment.

\section{A positive Laplacian does not itself supply a complex}

For finite inner-product spaces, let $D_0:C^0\to C^1$ and
$D_1:C^1\to C^2$ be arbitrary linear maps. Define
$L_1=D_0D_0^*+D_1^*D_1$. Positivity holds with no complex assumption:
\[
 \langle w,L_1w\rangle=\|D_0^*w\|^2+\|D_1w\|^2\ge0,
 \qquad \ker L_1=\ker D_0^*\cap\ker D_1.
\]
However, an orthogonal exact/coexact/harmonic decomposition requires
$D_1D_0=0$.

\begin{proposition}[Exact hypothesis for orthogonal Hodge components]
The subspaces $\operatorname{im}D_0$ and $\operatorname{im}D_1^*$ are
orthogonal if and only if $D_1D_0=0$. Under that hypothesis,
\[
 C^1=\operatorname{im}D_0\ \mathbin{\oplus^{\perp}}
 \operatorname{im}D_1^*\ \mathbin{\oplus^{\perp}}\ \ker L_1.
\]
\end{proposition}
\begin{proof}
For every $x,y$,
$\langle D_0x,D_1^*y\rangle=\langle D_1D_0x,y\rangle$.
Vanishing of all these pairings is equivalent to $D_1D_0=0$ by
nondegeneracy. The orthogonal complement of the sum of the two images is
$\ker D_0^*\cap\ker D_1=\ker L_1$. Finite dimensionality gives the
orthogonal direct sum. For arbitrary maps the two images can overlap;
positivity of $L_1$ does not remove that overlap.
\end{proof}

For covariant quiver differences, Theorem~\ref{thm:typed-transport-telescope}
identifies the curvature defect of each based face as $I-H$. An
approximately flat numerical connection is therefore not an exact chain
complex. Its residual norm can be reported and penalized, but its
``cohomology'' must not be treated as the cohomology of an exact differential
unless the relation is established. For gentle algebras the corresponding
path relations must likewise be checked in the actual representation.
Continuum ellipticity, mesh-uniform coercivity and convergence remain
additional analytic requirements.

\section{What a prefix-derived certificate can add}

The mutual-information identities in ToricGT II--IV need an explicit
distinction between the entire available prefix and a lossy model state.
The source definition makes its certificate measurable from the full
prefix, so it cannot add Bayes information to that same full prefix.

\begin{proposition}[Deterministic certificates have zero extra Bayes information]
Let $B,X,Z$ be finite random variables and suppose $Z=f(X)$. Then
\[
 H(B\mid X,Z)=H(B\mid X),\qquad I(B;Z\mid X)=0.
\]
The same conditional-expectation conclusion holds when $Z$ is measurable
with respect to a sigma-field $\mathcal F$ and $B$ has a finite alphabet.
\end{proposition}
\begin{proof}
For every $x$ of positive probability, the event $X=x$ already fixes
$Z=f(x)$. Thus the conditional distribution of $B$ given $(X,Z)$ is the
conditional distribution given $X$. Their conditional entropies coincide,
and their difference is the displayed mutual information. In the sigma-field
form, $\sigma(\mathcal F,Z)=\mathcal F$, so both conditional expectations
condition on the same sigma-field.
\end{proof}

This does not imply that computable certificates are useless. Let
$S=g(X)$ be the representation actually available to a bounded predictor.
In general $I(B;Z\mid S)$ can be positive, even though
$I(B;Z\mid X)=0$. For example, take $X$ uniform on $\{0,1,2,3\}$,
$B=Z=X\bmod2$, and $S=\lfloor X/2\rfloor$. Given $S$, $B$ is a fair bit;
given $(S,Z)$, it is known. The corresponding improvement is one bit, while
the improvement beyond complete $X$ is zero.

\begin{proposition}[Actual predictor improvement includes approximation error]
Let $q_0(\cdot\mid S)$ and $q_1(\cdot\mid S,Z)$ be normalized predictors
positive on the support of the target, with finite expected log losses
$L_0,L_1$. In base-two units,
\begin{align*}
 L_0-L_1={}&I(B;Z\mid S)\\
 &+\mathbb E\,D_{\rm KL}\bigl(p(\cdot\mid S)\|q_0(\cdot\mid S)\bigr)\\
 &-\mathbb E\,D_{\rm KL}\bigl(p(\cdot\mid S,Z)\|q_1(\cdot\mid S,Z)\bigr).
\end{align*}
\end{proposition}
\begin{proof}
For each conditioning value, insert and subtract $-\log_2p(B\mid\cdot)$
in the expected log loss. This gives conditional entropy plus conditional
KL divergence for each predictor. Subtract the two decompositions and use
the defining entropy difference for conditional mutual information.
\end{proof}

A deterministic topology, toric or proof certificate can therefore reduce
the approximation or optimization error of a restricted neural predictor.
That is a measurable training benefit, not new information beyond its
complete source data. A deployable certificate should be evaluated by
held-out predictive loss and its separately accounted artifact and runtime
cost. Leakage control still requires every certificate to be available
before the predicted observation; it does not require its usefulness to
come from a positive conditional mutual information given all that data.
This distinction also prevents a finite topological invariant from being
mistaken for extra evidence about a physical system.

\section{Damping and stochastic limits require their own hypotheses}

For a map $F$ with Lipschitz bound $L$ and $0<\eta\le1$, the relaxed map
$G_\eta=(1-\eta)I+\eta F$ has the upper bound
$\operatorname{Lip}(G_\eta)\le(1-\eta)+\eta L$.
When $L>1$, that upper bound is greater than one for every positive $\eta$.
It cannot certify contraction by decreasing $\eta$.

\begin{proposition}[Damping cannot repair every expanding cycle]
\label{prop:transport-damping-obstruction}
For $F(x)=2x$ on $\mathbb R$, $G_\eta(x)=(1+\eta)x$, so its exact
Lipschitz constant is greater than one for every $\eta>0$.
\end{proposition}
\begin{proof}
For distinct $x,y$, the ratio
$|G_\eta(x)-G_\eta(y)|/|x-y|$ equals $1+\eta>1$.
\end{proof}

The pinned manuscript's assertion that the displayed bound restores
contraction for $L>1$ is therefore incorrect. Damping can still stabilize
other maps: for $F(x)=-2x$, the relaxed multiplier is $1-3\eta$, whose
absolute value is less than one for $0<\eta<2/3$. The correct test concerns
the actual map, a valid norm bound, or a justified local Jacobian criterion;
an unnecessarily large upper bound is not an exact Lipschitz constant.

\begin{proposition}[A sufficient dissipativity condition for relaxation]
\label{prop:transport-dissipative-relaxation}
Let $b$ be a map on a real Hilbert space and suppose that, for all $x,y$,
\[
 \langle b(x)-b(y),x-y\rangle\le-\mu\|x-y\|^2,
 \qquad \|b(x)-b(y)\|\le L_b\|x-y\|,
\]
where $0<\mu\le L_b$. If $0<\eta<2\mu/L_b^2$, then
$x\mapsto x+\eta b(x)$ is a contraction with bound
$\sqrt{1-2\mu\eta+L_b^2\eta^2}<1$.
\end{proposition}
\begin{proof}
For $d=x-y$ and $e=b(x)-b(y)$, expand
$\|d+\eta e\|^2=\|d\|^2+2\eta\langle d,e\rangle+\eta^2\|e\|^2$.
The assumptions bound this by
$(1-2\mu\eta+L_b^2\eta^2)\|d\|^2$.
The coefficient is nonnegative since its minimum over $\eta$ is
$1-\mu^2/L_b^2\ge0$, and it is less than one on the stated interval.
Taking square roots proves the assertion.
\end{proof}

For $G_\eta=(1-\eta)I+\eta F$, apply the proposition to $b=F-I$.
This supplies a usable sufficient condition rather than an inference from
$\operatorname{Lip}(F)>1$. A learned skew-plus-dissipative fluid closure may
preserve the energy of one trajectory while failing this two-trajectory
condition when its coefficients depend on the state; that additional
derivative control must be established separately.

The stochastic ``embedding fluid'' passage needs a second qualification.
The continuous/hybrid GFlowNet theory of \citet{ToricAuditLahlou2023}
does not, by its balance identities alone, establish a small-step limit.
Let a small update have conditional mean
$\mathbb E[\Delta Z\mid Z=z]=b(z)\Delta t+o(\Delta t)$ and conditional
second moment
$\mathbb E[\Delta Z\Delta Z^T\mid Z=z]=a(z)\Delta t+o(\Delta t)$.
For a smooth compactly supported test function $\varphi$, Taylor expansion,
with a controlled remainder, gives
\[
 \frac{\mathbb E[\varphi(Z+\Delta Z)-\varphi(Z)\mid Z=z]}{\Delta t}
 \longrightarrow b\cdot\nabla\varphi+
 \frac12\sum_{i,j}a_{ij}\partial_i\partial_j\varphi.
\]
Under the additional tightness, moment and convergence assumptions needed
for a diffusion limit, the density consequently satisfies
\[
 \partial_t\rho=-\nabla\cdot(b\rho)
 +\frac12\sum_{i,j}\partial_i\partial_j(a_{ij}\rho).
\]
This conservative Markov limit has no independently generated source or sink.
Terms $\sigma_{\rm src}-\sigma_{\rm sink}$ require separately specified
branching, killing or injection dynamics, with their own scaling and
convergence hypotheses; they do not follow from the displayed conditional
first and second moments alone.
The pure continuity equation follows only when the second-order term
vanishes in the chosen scaling, or when the current has been defined to
include that diffusive contribution. A GFlowNet balance objective by itself
does not supply this limit. Likewise an antisymmetric-Jacobian penalty is a
useful proposed regularizer, not a theorem that semantic circular reasoning
has been eliminated or that physical Navier--Stokes dynamics is reproduced.

\section{Quantum and gentle-algebra versions of the transport contract}

The transport identity extends directly to linear maps between matrix
spaces. Let $\Phi_j$ be quantum channels and let $\rho_j$ be proposed
density operators. Set $r_j=\rho_j-\Phi_j(\rho_{j-1})$. The $r_j$ are
Hermitian trace-zero discrepancies, not density operators. Applying
Theorem~\ref{thm:typed-transport-telescope} to the underlying real vector
spaces of Hermitian matrices gives
\[
 \sum_j(\Phi_k\circ\cdots\circ\Phi_{j+1})(r_j)
 =\rho_k-(\Phi_k\circ\cdots\circ\Phi_1)(\rho_0).
\]
Trace-norm contractivity of each channel on Hermitian discrepancies then
bounds the final discrepancy by $\sum_j\|r_j\|_1$. This uses actual channel
assumptions; an arbitrary learned real connector does not acquire them by
being called quantum. A noisy cycle need not have identity holonomy, and
no inverse quantum operation is assumed.

For completeness, this contractivity has a short finite-dimensional proof.
For Hermitian $Y$, spectral decomposition gives
$\|Y\|_1=\max_{-I\le K\le I}\operatorname{Tr}(KY)$.
The adjoint of a positive trace-preserving map is positive and unital, so
$-I\le K\le I$ implies $-I\le\Phi^*(K)\le I$. Therefore
$\operatorname{Tr}(K\Phi(Y))=\operatorname{Tr}(\Phi^*(K)Y)\le\|Y\|_1$.
Taking the maximum proves $\|\Phi(Y)\|_1\le\|Y\|_1$; the triangle
inequality applied to the telescope gives the stated error bound. Complete
positivity supplies the stronger physical requirement of positivity after
adjoining an arbitrary ancillary system.

\begin{proposition}[Trace preservation obstructs a zero path relation]
\label{prop:transport-trace-zero-obstruction}
Suppose $\Phi$ and $\Psi$ are composable trace-preserving linear maps on
matrix spaces containing a normalized state $\rho$. Then
$\Psi\circ\Phi$ is not the zero map. Hence a nontrivial relation
$ba=0$ of a gentle path algebra cannot be represented by making both
composable arrows $a,b$ into trace-preserving channels on those full spaces.
\end{proposition}
\begin{proof}
Trace preservation gives
$\operatorname{Tr}(\Psi(\Phi(\rho)))=\operatorname{Tr}\rho=1$.
The zero map would instead produce trace zero. This contradiction uses
only linearity, normalization and trace preservation, so it applies in
particular to completely positive trace-preserving channels.
\end{proof}

The proposition is formalized, with an abstract linear trace functional, as
\begin{quote}\small
\path{TransportContracts.trace_preserving_composite_nonzero}.
\end{quote}
A legitimate quantum interpretation of
forbidden paths can use completely positive trace-non\-increasing instrument
branches, explicit rejection flags, or linear amplitude maps whose ranges
and kernels enforce the relation. For instrument branches, a forbidden
composition has zero probability; the complete instrument, including all
outcomes, is trace preserving. Normalizing a postselected branch is a
separate nonlinear state update and its success probability must be retained.
This distinction aligns the gentle algebra with operational quantum
information without assigning impossible channel semantics to its arrows.

\section{Tropical approximation is not an automatic universality transfer}

Finite ReLU networks give continuous piecewise-affine functions. With
arbitrary real slopes they can be expressed using differences of maxima of
real affine functions. Classical lattice-based toric geometry imposes
integrality or rationality data beyond this statement. Nonzero biases also
produce polyhedral complexes rather than necessarily conical fans. These
are compatible descriptions once the category and homogenization convention
are specified.

Softmax is smooth and is not a finite piecewise-affine map on an open
nontrivial logit region. A local piecewise-affine surrogate therefore needs
an approximation error bound. Moreover, dot-product logits with both query
and key variable are bilinear. A maximum of such scores need not be
piecewise affine in the joint inputs. The finite maximum-reduction identity
used by ring evaluation remains exact, independently of whether the score
functions are tropical polynomials.

The graph-universality argument in the Tropical Quivers manuscript also
requires a measure-specific step. The cited primary transformer theorem
is Yun et al.~\cite{ToricAuditYun2020}; the separate TokenGT construction is
\citet{ToricAuditTokenGT2022}. Yun et al.'s Theorems~2--3 use
$L^p$, $1\le p<\infty$, rather than the uniform norm asserted by the
fixed-cell theorem in the program's Paper~I.
Approximation in ambient Lebesgue
$L^p$ does not in general imply approximation on a lower-dimensional
tokenized graph family. For example,
$g_\epsilon(t)=\max(1-|t|/\epsilon,0)$ has
$\|g_\epsilon\|_{L^p(\mathbb R)}^p=2\epsilon/(p+1)$, tending to zero,
while $g_\epsilon(0)=1$ for every $\epsilon>0$. A tokenized domain supported
at that point sees no convergence. A valid transfer must specify the
graph-domain measure or prove uniform approximation there, and must match
the graph-relabeling action to the token symmetry group. Continuous
detokenization, finite padded sizes and admissible output types must also
be retained. Expressive approximation under those assumptions does not
establish that gradient descent finds the approximant or that unseen
physical distributions generalize.

The bundled transformer-geometry reference is version 1 of
\texttt{2604.14727}. Its current version 2 \cite{ToricAuditSuLiu2026}
explicitly distinguishes fixed-key
top-one routing, a log-lifted value parameterization, and finite-temperature
approximations.
That distinction should be retained when importing its geometry into a
trainable quiver with jointly varying keys and queries. The separate
tropical-attention paper \cite{ToricAuditHashemi2025} proposes a specified
algebraic kernel and reports algorithmic benchmarks. Its Appendix~C also
contains the localized errors repaired next, so this extension does not
import that proof package wholesale or transfer its claims to a different
dot-product-plus-maximum operation.

\section{Independent repairs to the tropical-attention reference}

Version~2 of \citet{ToricAuditHashemi2025}, Appendix~C's proof of
Lemma~3.1, uses $v_\lambda(a)=\log a-\lambda$ for $a>0$ and claims a
non-Archimedean additive inequality. This logarithm is an Archimedean
coordinate change. Neither a constant offset nor projective normalization
turns ordinary addition into exact maximum.

\begin{proposition}[Exact logarithmic identities and dequantization error]
\label{prop:tropical-log-repair}
For $a,b>0$ and real $\lambda$,
\[
 v_\lambda(ab)=v_\lambda(a)+v_\lambda(b)+\lambda.
\]
For any $m\ge1$, $z_1,\ldots,z_m\in\mathbb R$ and $\tau>0$,
\[
 M\le\tau\log\left(\sum_{j=1}^m e^{z_j/\tau}\right)
 \le M+\tau\log m,
 \qquad M=\max_j z_j.
\]
In particular, for $a=b=1$,
$v_\lambda(a+b)=\log2-\lambda>
\max\{v_\lambda(a),v_\lambda(b)\}$.
\end{proposition}
\begin{proof}
The product identity follows by expanding $\log(ab)=\log a+\log b$.
Factoring $e^{M/\tau}$ out of the sum leaves a number between one and
$m$. Apply the increasing logarithm and multiply by $\tau>0$ to obtain
the two bounds. The displayed counterexample follows from $\log2>0$.
\end{proof}

Thus the source's product offset has the wrong sign, and its additive
inequality omits a strictly positive term even on positive inputs.
The valid limiting statement is convergence to maximum as $\tau\downarrow0$
for a fixed finite number of terms, with the explicit error above. If the
number of terms varies, $\tau\log m$ must also be controlled.
For signed or complex terms, cancellation requires additional information,
as proved in Proposition~\ref{prop:reduction-signed-order}. These
corrections do not invalidate ordinary maximum-plus dynamic programming.

The same Appendix~C, Lemma~C.1, uses a large finite penalty as if it were
the exact value $-\infty$. A finite mask can be exact on a bounded input
domain, but it is not an input-independent annihilator on unbounded values.

\begin{proposition}[When a finite tropical mask is exact]
\label{prop:tropical-finite-mask}
Suppose a desired maximum has at least one admissible candidate of score
at least $L$. Suppose every excluded unpenalized candidate has score at
most $U$. Giving every excluded candidate an additive penalty $-M$ makes
it strictly unable to attain the maximum whenever $M>U-L$.
Without an upper bound on excluded values, no fixed finite $M$ guarantees
this conclusion.
\end{proposition}
\begin{proof}
Every penalized excluded score is at most $U-M<L$, below an admissible
score. Conversely, with an admissible score zero and any finite penalty
$M$, an excluded raw score $M+1$ becomes one and wins.
\end{proof}

An explicit $-\infty$ structural mask, or a proved range and sufficient
finite margin, repairs the construction. An all-masked query needs its own
defined output. Bounds on learned input features and their perturbations
must be included in a claimed finite-margin certificate.
The source also assigns $-\infty$ to excluded value entries in that proof.
If such entries are genuinely enforced, they remove those candidates
independently of the finite score penalty. The defect identified here is
the inference from a large finite penalty alone; it is not a counterexample
to a separately and exactly masked maximum.

\begin{proposition}[Exact-step and bounded-length tropical paths]
\label{prop:tropical-horizon-repair}
Write $\oplus=\max$ and $\otimes=+$ for scalar tropical operations, with
their induced matrix operations. Let $D$ be a finite tropical matrix and
let $I$ have zero diagonal and $-\infty$ elsewhere. Then
\[
 v_{t+1}=D\otimes v_t\quad\Longrightarrow\quad
 v_t=D^{\otimes t}\otimes v_0,
\]
whereas
\[
 w_{t+1}=w_t\oplus(D\otimes w_t),\quad w_0=v_0
 \quad\Longrightarrow\quad
 w_t=\left(\bigoplus_{j=0}^{t}D^{\otimes j}\right)\otimes v_0.
\]
Thus a pure $t$-step update and a finite Kleene closure are different
operations in general.
\end{proposition}
\begin{proof}
The first identity is induction and associativity. For the second,
distribute $D$ over the inductive sum; adding the previous sum retains
precisely powers zero through $t+1$. Idempotence removes duplicate terms.
For a direct counterexample take the one-by-one matrix $D=(-1)$ and
$v_0=0$. At $t=2$ the pure update is $-2$, while the maximum of the
zeroth, first and second powers applied to $v_0$ is zero. Even the source's
convention summing powers zero through $t-1$ gives zero, not $-2$.
\end{proof}

This repairs the mismatch in Theorem~C.5 of the reference. For maximum-plus
longest-path closure, deleting a repeated cycle can only improve a path
when every cycle has nonpositive total weight. Under that hypothesis a
best finite path can be chosen with at most $n-1$ edges, with the length-zero
path allowed. For minimum-plus shortest paths the corresponding hypothesis
is absence of negative-weight cycles. The sign and semiring cannot be
interchanged. On graphs with profitable cycles, a bounded-horizon
computation is still defined, but it is not the unbounded closure.

These are complete finite arguments for the operations actually used.
They do not assert that a trained attention stack has learned the exact
parameters, nor that any method solves arbitrary NP-hard instances in
polynomial resources. The end-to-end scientific contract remains a measured
prediction together with independently checked algebraic structure.

\section{A nonnormal obstruction to a Jacobian power certificate}

The Tropical Quivers manuscript's section on Jacobian spectral-radius
estimation iterates $v\mapsto Jv/\|Jv\|$ but labels the resulting estimate
as the largest singular value. For nonnormal matrices these are different
quantities. The following exact witness isolates the error without any
floating-point approximation.

\begin{proposition}[A nilpotent Jacobian with arbitrarily large gain]
\label{prop:nonnormal-jvp}
For any real $M>1$, put
\[
 J_M=\begin{pmatrix}0&M\\0&0\end{pmatrix},\qquad
 e_2=\begin{pmatrix}0\\1\end{pmatrix}.
\]
Then $J_M^2=0$, $\rho(J_M)=0$, and $\|J_M\|_2=M$.
Nevertheless $\|J_Me_2\|_2^2=M^2>1$, while the next application of
$J_M$ is zero. In particular, a vanishing or small eigenvalue power probe
does not certify one-step Euclidean contraction.
\end{proposition}
\begin{proof}
Direct multiplication gives $J_M^2=0$ and
$J_M^TJ_M=\operatorname{diag}(0,M^2)$, whereas
$J_MJ_M^T=\operatorname{diag}(M^2,0)$. Thus $J_M$ is nonnormal and
both eigenvalues are zero. For $x=(x_1,x_2)$,
$\|J_Mx\|_2^2=M^2x_2^2\le M^2\|x\|_2^2$, with equality at
$x=e_2$. This proves the norm formula and all remaining assertions.
\end{proof}

For a general real Jacobian, singular-norm power iteration uses
$J^TJ$, requiring both a Jacobian-vector product and its adjoint product.
Even that correction does not turn a finite numerical estimate into a
certified upper bound: a Rayleigh quotient samples a direction and can
miss the maximizing singular direction. A contraction certificate needs a
proved upper bound, a validated enclosure, or an explicitly controlled
operator parametrization. The nilpotent example can have powers that
eventually vanish while still amplifying perturbations at the first step.

The companion formalizes its explicit matrix products, nonnormality and
squared coordinate gain. The spectral-radius and operator-norm conclusions
are the written finite-dimensional argument above; the corresponding Lean
components are not represented as a complete analytic norm formalization.

\section{Path multiplicity separates tropical growth from operator stability}

The max-plus cocycle in \citet{ToricAuditQuivers2026} records the largest
logarithmic path weight. An ordinary matrix product adds contributions
from many paths. Its operator growth therefore needs a multiplicity term.
This is a qualification of the source's contraction interpretation, not
a rejection of its max-plus subadditivity calculation.

Let $A_t:\mathbb C^{d_{t-1}}\to\mathbb C^{d_t}$, $1\le t\le T$ and $T\ge1$, with
each $d_t\ge1$, and set $P_T=A_T\cdots A_1$. An admissible path from $j$
to $i$ is a sequence $(i_0,\ldots,i_T)$ with $i_0=j$, $i_T=i$ and
all $A_t(i_t,i_{t-1})\ne0$. Let $C_T(i,j)$ count these paths and let
$W_T(i,j)$ be their maximum absolute product weight, taking both values
to be zero when no admissible path exists. Thus $W_T$ is a maximum-product
matrix product; its logarithm, with $\log0=-\infty$, is the corresponding
max-plus product. Let $S_t$ be the zero--one support matrix of $A_t$.

\begin{theorem}[Finite path-multiplicity bound]
\label{thm:tropical-path-multiplicity}
With the preceding notation,
\[
 |P_T(i,j)|\le (|A_T|\cdots|A_1|)(i,j)
 \le C_T(i,j)W_T(i,j),
 \qquad C_T=S_T\cdots S_1.
\]
Writing $N_T=\max_i\sum_jC_T(i,j)$ and
$m_T=\max_{i,j}W_T(i,j)$ gives the induced row-sum norm bound
\[
 \|P_T\|_\infty
 \le\max_i\sum_j C_T(i,j)W_T(i,j)
 \le N_Tm_T.
\]
For square $n$-dimensional arrows, $N_T\le n^T$. Consequently, whenever
$m_T>0$ eventually and
$\limsup_{T\to\infty}T^{-1}\log m_T+\log n<0$, the ordinary matrix
products decay exponentially in this norm.
\end{theorem}
\begin{proof}
Expand each entry of $P_T$ as the sum over the intermediate indices.
The triangle inequality bounds its absolute value by the sum of the
absolute path weights. Each nonzero summand is at most $W_T(i,j)$,
and exactly $C_T(i,j)$ such paths occur. Multiplying the support matrices
counts those paths, proving the entrywise assertions. Sum over $j$ and
maximize over $i$ to obtain the norm inequalities. In dimension $n$, a
fixed final index leaves $T$ freely chosen earlier indices, so there
are at most $n^T$ paths ending at that index. The strict limsup inequality
then gives constants $c>0$ and $T_0$ with
$n^Tm_T\le e^{-cT}$ for $T\ge T_0$.
\end{proof}

If $m_T=0$ at any time, every path of that length vanishes, so all later
products are zero and the conclusion holds directly. More precise bounds
can retain $N_T$ or the entire entrywise table instead of $n^T$.
Using $\log(\epsilon+|A_t(i,j)|)$ with $\epsilon>0$, as in the source,
increases path weights and produces a conservative maximum-product upper
bound; it does not remove the number of summed paths. Conversely, signed
or complex cancellation can make $P_T$ smaller than the positive bound.
These inequalities do not recover such cancellation from magnitudes.

\begin{proposition}[Decaying maximal paths with an expanding linear map]
\label{prop:tropical-branching-counterexample}
For $n\ge2$ and $1/n<a<1$, let $A=a\mathbf1\mathbf1^T$ on
$\mathbb R^n$. Every length-$k$ path has weight $a^k$, so its tropical
growth exponent is $\log a<0$. Nevertheless
\[
 A\mathbf1=na\mathbf1,\qquad
 A^k=n^{k-1}a^k\mathbf1\mathbf1^T,\qquad
 \|A^k\|_\infty=(na)^k>1 \quad(k\ge1).
\]
The path-multiplicity norm bound is attained.
\end{proposition}
\begin{proof}
Each entry is $a$, and each row has $n$ entries. The eigenvector identity
follows by row summation. Since
$(\mathbf1\mathbf1^T)^2=n\mathbf1\mathbf1^T$, induction gives the
power formula. Its row sum is $(na)^k$. There are $n^{k-1}$ paths between
any prescribed pair of endpoints and $n^k$ paths ending at a prescribed
final vertex, so both multiplicity inequalities are equalities.
\end{proof}

For example, $n=2$, $a=3/4$ gives individual path decay $(3/4)^T$ and
ordinary growth $(3/2)^T$. Thus a negative maximal-path exponent alone
cannot certify contraction of a learned fluid closure, a quiver update,
or a quantum-amplitude surrogate. A separately proved norm inequality
or a preserved stochastic/channel structure is required. The exact
regression computes support counts and maximum weights by dynamic
programming and compares them against exhaustive path enumeration on
small, possibly rectangular, signed rational examples. It is a validation
of finite products, not an estimate of an infinite-time physical exponent.

\section{Weighted small-gain certification at a fixed rate}

The weighted small-gain criterion in \citet{ToricAuditQuivers2026} is
useful, but its instruction to minimize the weighted row bound by a
linear program needs a qualification. Suppose $G\ge0$ is a finite
matrix of valid block Lipschitz gains. For positive weights $w_i$ define
$\|x\|_{\infty,w}=\max_i\|x_i\|/w_i$. Its row bound is
\[
 q(w)=\max_i \frac{(Gw)_i}{w_i}.
\]

\begin{proposition}[Fixed-rate feasibility is linear]
\label{prop:weighted-gain-feasibility}
For fixed $q\ge0$, $q(w)\le q$ is equivalent to the componentwise
linear inequalities $Gw\le qw$, $w_i>0$. If fixed bounds
$0<\ell_i\le w_i\le u_i$ are imposed, feasibility is a linear-program
problem. Feasibility is monotone in $q$. Jointly minimizing $q$ over
$w$ in this representation is not a linear program because $qw_i$
contains products of variables.
\end{proposition}
\begin{proof}
Multiply each row-ratio inequality by its positive denominator $w_i$.
At fixed $q$, every resulting coefficient is fixed, so the inequalities
and box bounds are linear. If a weight vector is feasible at $q$, it
remains feasible at every larger rate because its coordinates are
positive. If $q$ is also variable, $qw_i$ is bilinear, which establishes
the final qualification.
\end{proof}

One may bisect a certified bracket of rates and solve the fixed-rate
feasibility problems. A returned numerical witness should still be
checked with exact arithmetic or validated error bounds before it is used
as a proof of contraction. The checker supplied here verifies a given
rational witness and does not claim to solve the optimization problem.
For
\[
 G=\begin{pmatrix}0&2\\1/8&0\end{pmatrix},\qquad
 w=\begin{pmatrix}4\\1\end{pmatrix},
 \qquad Gw=\tfrac12w,
\]
the unweighted row bound is two, while the weighted bound is $1/2$.
For any positive feasible weights, the two inequalities multiply to
$1/4\le q^2$, so a proposed rate $q=2/5$ is impossible. These exact
equalities and the rejection are included in the regression. Bounds on
$G$ must hold over the claimed domain; a few sampled Jacobians do not
establish them.

\section{Proper scoring must include its probability boundary}

The cross-entropy theorem in \citet{ToricAuditOracle2026} optimizes over
$q(Z)\in(0,1)$ but does not exclude deterministic conditional labels.
When $\eta(Z)=\mathbb P(Y=1\mid Z)$ equals zero or one, its stated
minimizer lies outside that open interval.

\begin{proposition}[Binary log loss with correct endpoint conventions]
\label{prop:oracle-logloss-endpoints}
For $\eta\in[0,1]$, extend
\[
 R_\eta(q)=-\eta\log_2q-(1-\eta)\log_2(1-q)
\]
to $q\in[0,1]$ by assigning zero to a zero-probability term and
$+\infty$ to a positive-probability event assigned probability zero.
Then $R_\eta$ is uniquely minimized at $q=\eta$. Over the restricted
domain $q\in(0,1)$ the minimizer exists precisely when $0<\eta<1$;
for $\eta=0$ or $1$, its infimum is zero and is not attained.
\end{proposition}
\begin{proof}
If $0<\eta<1$, both endpoints have infinite risk. In the interior,
$R_\eta'(q)=(q-\eta)/(q(1-q)\ln2)$, which is negative below
$\eta$ and positive above it. Thus the interior minimum is unique.
If $\eta=0$, the risk is $-\log_2(1-q)$, uniquely minimized at zero.
Every interior $q$ has positive risk and is improved by replacing it
with $q/2$; the risk tends to zero as $q\downarrow0$. If $\eta=1$,
the risk is $-\log_2q$, uniquely minimized at one. Every interior
$q$ is improved by $(1+q)/2$, and the risk tends to zero as
$q\uparrow1$. These cases prove all assertions.
\end{proof}

Applying this argument conditionally gives the Bayes statement for
measurable forecasts with the corresponding admissible range. A sigmoid
with finite logits produces only interior probabilities, so deterministic
targets need not admit a finite-logit optimum. The proper-scoring theorem
does not by itself prove empirical calibration of a restricted, finitely
trained oracle. Calibration and useful held-out control remain measured
properties, with the prefix and artifact accounting established earlier.

\section{Predictive synergy is not determined by geometric overlaps}

The last context-utility theorem of \citet{ToricAuditMemory2026} proposes
an additive information decomposition over independent sheaf components,
with overlap corrections determined by Cech data. The usual chain rule
for mutual information contains conditional statistical terms; independence
of source components alone does not remove predictive synergy.

\begin{proposition}[Independent disjoint components can carry one bit jointly]
\label{prop:context-xor-synergy}
Let $X,Y$ be independent fair bits, let $B=X\mathbin{\mathrm{xor}}Y$,
and let the additive statistic be $T=X+Y$. Then $B$ is determined by
$T$, but
\[
 I(B;X)=I(B;Y)=0,\qquad
 I(B;T)=I(B;X,Y)=1.
\]
The two source components may be stored on disjoint pieces with no
overlap restriction data.
\end{proposition}
\begin{proof}
The four pairs $(X,Y)$ have probability $1/4$. Conditional on either
value of either source bit, $B$ remains a fair bit, so both individual
mutual informations vanish. The target is a fair bit overall and equals
$T\bmod2$, so its conditional entropy given $T$ is zero. It is also
determined by the pair $(X,Y)$. Its unconditional entropy is one bit,
which proves the other equalities. The storage locations of the two
source bits impose no restriction on this joint probability law.
\end{proof}

For arbitrary finite variables and conditioning context $C$, the exact
decomposition is
\begin{align*}
 I(B;X,Y\mid C)
 &=I(B;X\mid C)+I(B;Y\mid X,C)\\
 &=I(B;X\mid C)+I(B;Y\mid C)
   +I(X;Y\mid B,C)-I(X;Y\mid C).
\end{align*}
Both identities follow by expanding conditional mutual information as
differences of conditional entropies and applying the entropy chain rule.
Even if $I(X;Y\mid C)=0$, conditioning on the target can make the
sources dependent. In the example, $I(X;Y\mid B)=1$, and there is
no source overlap whose absence would force this correction to vanish.
Additional probabilistic factorization assumptions are needed to replace
these terms by a particular geometric-overlap model. Compatibility of local
sections alone is insufficient. If $X,Y$ are already determined by the
complete conditioning prefix, all extra conditional information is zero,
as proved earlier; the example applies when the actual predictor has
retained only a smaller context.

The regression checks the complete rational probability table, all
pairwise-independence equalities, deterministic recovery from $T$, and
the conditional dependence given $B$. Its reported one-bit values follow
from these fair and deterministic distributions. The general entropy and
log-loss arguments in this section are written proofs; they are not added
to the count of Lean-checked declarations.

\section{Implications for the research program}

The repaired transport identity provides a shared finite contract for
surface cochains, quiver connectors, quantum instruments and learned
trajectory residuals. It can support a structure-preserving fluid model,
but a Navier--Stokes proof still needs the external construction's analytic
estimates. The corresponding GRH/Hilbert--P\'olya program still requires a
precise arithmetic spectral identification and operator-domain control;
finite self-adjoint matrices do not identify their spectra with zeta zeros.
A positive finite Hodge spectrum likewise does not establish a uniform
continuum Yang--Mills mass gap. Finally, a path-algebra certificate can
authenticate a structural computation only when an actual cryptographic
protocol binds it to a message and a threat model; difficult-looking
tropical equations are not, by themselves, hardness assumptions.

The exact regression script tests 288 generic rational paths, additional
rectangular paths, 128 flat cycles with noncommuting arrows, the prefix
counterexample, positive damping examples and the finite information
example. It also compares tropical dynamic programming against exhaustive
bounded-path enumeration and records the finite-mask and horizon
counterexamples. The new path-multiplicity checks also compare exact
products with an independent enumeration, including changing dimensions,
and retain a dense example attaining the bound. Fixed-rate small-gain
witnesses, the proper-scoring endpoint inequalities and the full XOR
probability table are checked separately. The logarithmic statements are ordinary mathematical
proofs above; the exact regression encodes their positive-input
counterexamples using the monotonicity of the logarithm, without treating
floating-point logarithms as formal certificates. The Lean log records
the exact theorem types and axiom dependencies for its seventeen selected
declarations. Those records state what has been checked and leave the
larger scientific and analytic obligations visible.
